\documentclass[11pt]{article}
\usepackage{amsmath,amssymb,amsthm}
\usepackage[margin=1.1in]{geometry}
\newtheorem{theorem}{Theorem}
\newtheorem{proposition}{Proposition}
\newtheorem{definition}{Definition}
\theoremstyle{remark}
\newtheorem{remark}{Remark}

\title{Recovering the primes in a dyadic interval from exact moments:\\
two constructive routes and an obstruction theorem}
\author{Carl Gribble\thanks{Draft v0.1, third in a series, companion to \emph{Prime counts and prime
sums from multiplication tables: a self-similar formula system}~\cite{CompanionI} and \emph{Certified
exact sums of primes over dyadic intervals via a ladder of generalized factorials}~\cite{CompanionII}
(same author). Developed in collaboration with Claude (Anthropic); all numerical claims were
machine-verified as described in Section~6. Responsibility for the content rests with the author. A full disclosure of AI use appears at the end of the paper.}}
\date{August 2026 --- draft}

\begin{document}
\maketitle

\begin{abstract}
The companion papers supply two exact ``moment channels'' for the primes in a dyadic window
$W = (m, 2m]$: integer power sums $\sum_{p \in W} p^k$, produced from divisor-lattice data with no
prime input at any stage \cite{CompanionI}, and weighted logarithmic moments
$\sum_{p \in W} Q(p)\log p$ for polynomial $Q$, produced from the primes below the window
\cite{CompanionII}. This note closes the circle: the channels determine the window primes
\emph{individually}. Two constructive decoders are given. From the $r{+}1$ integer moments of weights
$0, \dots, r$ ($r$ the number of window primes), Newton's identities yield the elementary symmetric
functions --- integers --- and the primes are read off as the integer roots of the resulting monic
polynomial; the entire pipeline contains no prime input and no primality test. Alternatively, a
\emph{single} real moment, Chebyshev's 1852 functional $\theta(2m) - \theta(m)$, known to
$\lceil \log_2 N \rceil + 2$ bits where $N = \prod_{p \in W} p \le 4^m$, determines $N$ by rounding,
and the window primes are exactly the divisors of $N$ in $W$, found by trial division; a Prony decoder
using $2r$ moments at intermediate precision interpolates between the two. Against these, an
obstruction: any scheme whose decoder receives at most $D$ bits determined by the configuration, and
which correctly decodes \emph{every} $r$-subset of the window's integers, has
$D \ge \log_2 \binom{m}{r} - 1 \asymp m \log\log m / \log m$. Hence no uniform scheme of bounded
description decodes the dyadic windows for all $m$, while the companion's certified determination of
the single aggregate $S(m)$ from $O(\log m)$ moments sits near the same floor read at entropy
$O(\log m)$: one pigeonhole, two entropy scales. The model restriction --- decoders blind to number
theory except through the delivered bits --- is discussed and shown to be necessary; Mills- and
Willans-type formulas live outside it, as does the sieve. All constructions are machine-verified: at
$m = 500$ both routes recover the $73$ primes of $(500, 1000]$ exactly, the second from a single
$695$-bit real number.
\end{abstract}

\section{Introduction}

The first companion paper \cite{CompanionI} develops Linnik's identity \cite{Linnik}, in the
computational tradition opened by McKenzie \cite{McKenzie}, into a formula system $S_j$ returning the
exact integers $\sum_{p \le y} p^{\,j}$ for every weight $j \ge 0$ from divisor-lattice data alone: no
prime, primality test, or sieve appears at any depth. The second \cite{CompanionII} constructs
generalized factorials $\prod_{i \le 2m} i^{\,Q(i)}$ with polynomial exponents whose valuations at
window primes are $Q(p)$, so that each polynomial weight yields one exact functional
$\sum_{m < p \le 2m} Q(p) \log p$, computable from the primes $\le m$; finitely many such moments,
closed by a certified sandwich and the integrality of the target, determine the aggregate
$S(m) = \sum_{m < p \le 2m} p$ exactly.

Both constructions output \emph{aggregates}. It is natural to ask how much of the window the moment
channels actually know. The answer, made precise here, is: everything. Finitely many moments determine
the window primes individually, by elementary and fully constructive decoders (Theorems~\ref{thm:int}
and \ref{thm:prod}, Remark~\ref{rem:prony}), and the cost of this determination --- counted in
delivered bits --- can be compared against an unconditional pigeonhole floor
(Theorem~\ref{thm:lb}). The comparison quantifies, in a restricted but honestly delimited model, the
folklore assertion that no short formula produces the primes: the floor grows like
$m \log\log m / \log m$ for the full window, while a single aggregate of entropy $O(\log m)$ admits the
companion's $O(\log m)$-moment certified determination. The impossibility of a bounded-length exact
scheme for the primes themselves and the possibility of a log-length certified procedure for one sum
are the same counting argument read at two entropy scales.

We state claims and non-claims plainly. Claimed: (i) the exact recovery theorems and their machine
verification; (ii) the bit-budget lower bound, with an explicit discussion of the model restriction
that makes it true; (iii) the accounting that unifies the three regimes (one exact real moment; many
finite-precision moments; too few bits). Not claimed: computational advantage --- both decoders cost
far more than sieving the window, by design, since the object of study is the information channel, not
the running time; and no unrestricted impossibility --- the primes are computable from $m$, so any
lower bound must, and ours does, constrain the decoder's access rather than the universe.

\section{The two moment channels}

Fix $m \ge 8$, write $W = (m, 2m]$, let $P(W)$ be the set of primes in $W$ and $r = |P(W)|$ (so
$r \ge 1$ by Bertrand's postulate). The channels of \cite{CompanionI, CompanionII} deliver:
\begin{align}
\text{(integer channel)} \qquad & s_k \;=\; \sum_{p \in P(W)} p^{\,k} \;=\; S_k(2m) - S_k(m)
\;\in\; \mathbb{Z}, \qquad k \ge 0,\label{eq:intch}\\
\text{(weighted channel)} \qquad & M_Q \;=\; \sum_{p \in P(W)} Q(p)\, \log p, \qquad Q \in \mathbb{Z}[t],
\label{eq:wch}
\end{align}
the first exactly, in integer arithmetic, with no prime input at any stage \cite{CompanionI}; the
second to any prescribed precision, from the primes $\le m$ together with logarithms of the integers
$\le 2m$ \cite[Prop.~1]{CompanionII}. The case $Q = 1$ of \eqref{eq:wch} is Chebyshev's identity
\cite{Cheb} in the form
$\theta(2m) - \theta(m) = \log\,(2m)! \,-\, \sum_{p \le m} v_p\big((2m)!\big)\log p$, used below. Only the two displayed facts \eqref{eq:intch} and \eqref{eq:wch}, with the stated
computability provenance, are imported from the companions; every proof in this paper is
self-contained given those displays.

\section{Route I: integer moments, Newton's identities, integer roots}

\begin{theorem}[Prime-free individual recovery]\label{thm:int}
Let $r = s_0 = S_0(2m) - S_0(m)$ and let $s_1, \dots, s_r$ be as in \eqref{eq:intch}. Define $e_0 = 1$
and, recursively,
\[
e_k \;=\; \frac{1}{k} \sum_{i=1}^{k} (-1)^{i-1} e_{k-i}\, s_i, \qquad 1 \le k \le r .
\]
Then every $e_k$ is an integer,
\[
f(x) \;=\; \sum_{k=0}^{r} (-1)^k e_k\, x^{\,r-k} \;=\; \prod_{p \in P(W)} (x - p),
\]
and $P(W) = \{\, c \in W \cap \mathbb{Z} : f(c) = 0 \,\}$. The entire pipeline --- moments, recursion,
and root reading by evaluation at the $m$ integers of $W$ --- consumes no prime and performs no
primality test; the primes emerge as the integer roots.
\end{theorem}

\begin{proof}
The $s_k$ are the power sums of the finite multiset $P(W)$ of distinct integers. Newton's identities
(see, e.g., \cite[Ch.~I, \S2]{Macdonald}) express the elementary symmetric functions $e_k$ of that
multiset by exactly the stated recursion; the $e_k$ are symmetric polynomials with integer
coefficients in integers, hence integers, and the recursion's divisions by $k$ are exact. The monic
polynomial with those symmetric functions is $\prod_{p}(x - p)$. A monic integer polynomial of degree
$r$ has at most $r$ integer roots; the $r$ window primes are roots; so the zero set of $f$ in $W$ is
exactly $P(W)$.
\end{proof}

The moment count is the smallest for which this decoder is defined: $r$ unknowns, $r$ power sums, plus
$s_0$ to know $r$. The price is paid in height rather than count: $\log_2 e_k$ grows like
$k \log_2 m$, so the total delivered data is $\Theta(r^2 \log m)$ bits (Section~7), and the encoder
must run the companion system at weights up to $r$.

\section{Route II: one real moment, rounding, trial division}

\begin{theorem}[Single-moment recovery]\label{thm:prod}
Let $N = \prod_{p \in P(W)} p$, so $N \mid \binom{2m}{m}$ and $N \le 4^m$. Put
$B = \lceil \log_2 N \rceil + 2$ (a priori, $B = 2m + 2$ suffices). Let $\widetilde{M}$ be any real
number with
\[
\bigl|\widetilde{M} - \big(\theta(2m) - \theta(m)\big)\bigr| \;\le\; 2^{-B},
\]
for instance the output of the extraction identity of \cite[Prop.~1]{CompanionII} at rung $0$,
evaluated in interval arithmetic at $B + O(1)$ bits from the primes $\le m$ and the logarithms of the
integers $\le 2m$. Then $N$ is the unique integer within distance $\tfrac12$ of
$\exp(\widetilde{M})$, and
\[
P(W) \;=\; \{\, c \in W \cap \mathbb{Z} \,:\, c \mid N \,\},
\]
recovered by trial division of $N$ by the $m$ integers of the window. No integer in $W$ is tested for
primality.
\end{theorem}

\begin{proof}
$\exp(\theta(2m) - \theta(m)) = N$ exactly, and for $|\delta| \le 2^{-B}$,
$|N e^{\delta} - N| \le N\, 2^{-B+1} < 2^{\lceil \log_2 N\rceil + 1 - B} = \tfrac12$, so rounding
recovers $N$; the divisibility $N \mid \binom{2m}{m}$ and the bound $N \le 4^m$ are the central
binomial facts underlying \cite{Cheb,Erdos}. For the divisor characterization: every $p \in P(W)$ divides
$N$ and lies in $W$; conversely a composite $c \in W$ has a prime factor $q \le \sqrt{2m} \le m$
(valid for $m \ge 2$), and $q$ does not divide $N$ since every prime factor of $N$ exceeds $m$; hence
$c \nmid N$.
\end{proof}

Theorem~\ref{thm:prod} is the ``internal Willans'' of the moment model: a single real number carries
the entire window. Two comments calibrate it. First, the information is priced in precision --- the
one moment must be delivered to $\approx \log_2 N \approx (\theta(2m)-\theta(m))/\log 2$ bits, and
Section~7 shows this is within a factor $\sim \log m / \log\log m$ of the unconditional floor, making
rung $0$ at high precision the most bit-efficient constructive decoder we know. Second, information
and extraction separate: the \emph{exact} real $\theta(2m)-\theta(m)$ determines $P(W)$ outright, but
generic extraction from $N$ is integer factorization of an $r\log m$-bit number; the window structure
rescues us, since trial division over $W$ is \emph{complete} --- composites cannot divide $N$ --- so
no factoring and no primality testing occur.

\begin{remark}[Prony interpolation]\label{rem:prony}
The weighted channel \eqref{eq:wch} with $2r$ moments recovers $P(W)$ at intermediate precision. With
$\sigma_k = \sum_p (\log p)\, p^k$ for $0 \le k \le 2r-1$, the Hankel matrix
$(\sigma_{i+j})_{0 \le i, j \le r-1} = V^{\mathsf T} D V$ is positive definite ($V$ the Vandermonde
matrix at the distinct primes, $D = \mathrm{diag}(\log p) \succ 0$), and the Prony system
$\sum_{j=0}^{r-1} c_j \sigma_{i+j} = -\sigma_{r+i}$ returns the coefficients of
$\prod_p (x - p) \in \mathbb{Z}[x]$, which rounding makes exact at sufficient precision; the roots are
then read off as in Theorem~\ref{thm:int}. Explicit precision bounds follow from Vandermonde-inverse
estimates using the gap $|p_i - p_j| \ge 2|i - j|$ between distinct odd primes, but are tedious; we
record instead the machine datum that at $m = 50$ a working precision of about $160$ bits suffices
(Section~6). Prony's memoir is \cite{Prony}.
\end{remark}

\section{The obstruction theorem}

The primes up to $2m$ are computable from the single input $m$; in the unrestricted world their
description length is $O(\log m)$ and no lower bound is possible. The meaningful question, and the
rigorous kernel of the no-closed-form folklore, concerns schemes whose decoders learn about the window
\emph{only through delivered data}. We formalize the channel and count.

\begin{definition}[Uniform scheme]\label{def:scheme}
Fix $m$ and $1 \le r \le m$. A \emph{$D$-bit uniform scheme} for $(m, r)$ is a pair $(\Phi, \Psi)$
where $\Phi$ maps each $r$-element subset $S \subseteq W \cap \mathbb{Z}$ to a bit string of length at
most $D$, and $\Psi(\Phi(S)) = S$ for every such $S$. No constraint is placed on the functions
themselves: $\Phi$ may be any quantization of any moments, and $\Psi$ any decoder --- but $\Psi$ sees
only the bits.
\end{definition}

\begin{theorem}[Bit floor]\label{thm:lb}
Every $D$-bit uniform scheme for $(m, r)$ satisfies $D \ge \log_2 \binom{m}{r} - 1$. In particular,
for $r = \pi(2m) - \pi(m) \sim m/\log m$,
\[
D \;\ge\; (1 + o(1))\, \frac{m \log_2(e \log m)}{\log m} \;\longrightarrow\; \infty ,
\]
so no family of uniform schemes of bounded description length decodes the dyadic windows for all $m$.
\end{theorem}

\begin{proof}
$\Psi \circ \Phi = \mathrm{id}$ forces $\Phi$ to be injective on the $\binom{m}{r}$ configurations;
there are $2^{D+1} - 1$ bit strings of length at most $D$; so $2^{D+1} > \binom{m}{r}$. The
asymptotic is $\log_2\binom{m}{r} = r \log_2(em/r) + O(\log m)$ with $r \sim m/\log m$.
\end{proof}

Three glosses carry the theorem's meaning. \emph{(a) The ensemble is the model.} The actual window
primes are one configuration per $m$; correctness over all $r$-subsets is what excludes decoders with
number theory built in. Mills's prime-representing function \cite{Mills} is a decoder whose constant
stores the answers; Willans's formula \cite{Willans} passes unbounded information through the
evaluation of factorials under a floor; a sieve is a decoder that ignores the channel and recomputes
from $m$. All three are outside Definition~\ref{def:scheme}, and their exclusion is precisely what
makes the folklore well-posed rather than false. The decoders of Theorems~\ref{thm:int} and
\ref{thm:prod} --- Newton's recursion, rounding, trial division --- use no property of the primes
beyond what the bits assert, are correct on arbitrary configurations of distinct window integers (for
Theorem~\ref{thm:prod}, arbitrary sets of primes in $W$; for Theorem~\ref{thm:int}, arbitrary distinct
integers), and hence live inside the model, where their budgets can be compared to the floor.
\emph{(b) Precision must be charged.} With exact reals even one moment suffices
(Theorem~\ref{thm:prod}), so a count-only lower bound is impossible; the bit budget $D$ is the honest
currency, and the count--precision \emph{tradeoff} is the object of study. \emph{(c) Information is
not extraction.} The floor prices what the bits can distinguish, not the work of deciding; the time
cost of prime determination is a separate and harder question, on which the elementary
$\sqrt{x}$-barrier literature \cite{Polymath} is the state of the art.

\section{Machine verification}

A reference implementation (\texttt{prime\_recovery.py} in the code repository \cite{Code}) exercises
both routes end to end, re-using the companion systems as published. At $m = 500$, so
$W = (500, 1000]$: the integer channel returned $r = s_0 = 73$ and $s_1, \dots, s_{73}$ with no prime
input at any stage (the companion system's exactness assertion --- every harmonic combination landing
on denominator $1$ --- held at all weights up to $73$); Newton's recursion divided exactly at every
step; the $73$ integer roots of $f$ in $W$ coincide with an independent sieve; and
$e_{73} = N$, a $695$-bit integer, all in under a second. Independently, rung $0$ evaluated by the
extraction identity at $1300$-bit working precision, from the primes $\le 500$ and the logarithms of
the integers $\le 1000$, satisfied $|\exp(\widetilde M) - N| \approx 4 \times 10^{-178}$; rounding
recovered the same $N$ (the cross-route identity $e_{73} = N$ is itself a check that the two channels
agree), and trial division over the window returned exactly the $73$ primes, with no primality test
performed. At $m = 50$, the Prony decoder of Remark~\ref{rem:prony} with $2r = 20$ weighted moments
recovered the ten primes of $(50, 100]$; a precision scan located the working threshold near $160$
bits, every rounded coefficient landing within $1/4$ of the integer coefficients of
$\prod_p (x - p)$. Trust base: big-integer arithmetic, the mpmath library, and Newton's identities;
nothing else is assumed.

\section{The accounting, the gap, and open questions}

Theorem~\ref{thm:lb} supplies the floor; the constructions supply budgets. At the verified scales:

\begin{center}
\begin{tabular}{lcccc}
\hline
scheme & window & moments & total delivered bits & $\log_2 \binom{m}{r}$ \\
\hline
Route I (Thm.~\ref{thm:int}) & $(500, 1000]$ & $74$ & $27{,}128$ & $\approx 296$ \\
Route II (Thm.~\ref{thm:prod})$^{\dagger}$ & $(500, 1000]$ & $1$ & $697$ & $\approx 296$ \\
Prony (Rem.~\ref{rem:prony})$^{\dagger}$ & $(50, 100]$ & $20$ & $\approx 3{,}200$ & $\approx 33$ \\
\hline
\end{tabular}
\end{center}

\noindent
Route I is uniform over arbitrary distinct-integer configurations and so lies fully inside
Definition~\ref{def:scheme}; its $\times 90$ overpay is the redundancy of the symmetric-function
encoding, $\Theta(r^2 \log m)$ bits against a floor of $\Theta(r \log\log m)$ --- the price of the
minimal moment \emph{count}. The daggered schemes are uniform over configurations of window
\emph{primes} only: over arbitrary integer configurations a single moment cannot even separate ---
$\{54, 96\}$ and $\{64, 81\}$ share their product --- so the injectivity of rung $0$ on prime
configurations is itself a unique-factorization fact, and the $697$-against-$296$ comparison bounds
the overpay of the product encoding \emph{granted} the side information that the configuration is
prime. That overpay has a clean shape: the product spends $\log_2 p \approx \log_2 m$ bits per prime
where the floor charges only $\log_2(em/r) \approx \log_2(e \log m)$, a ratio
$\log m / \log(e\log m)$, about $2.4$ at $m = 500$ and slowly divergent. Rung $0$ at high precision is
thus the most bit-efficient constructive decoder we know, and it is not optimal.

The unification promised in the introduction is now one sentence. The companion sandwich
\cite{CompanionII} determines the single aggregate $S(m)$ --- entropy $O(\log m)$ --- from
$K + 1 \asymp \log m$ moments of bounded precision, a polylogarithmic budget sitting near the floor at
the small entropy scale; Theorem~\ref{thm:lb} shows the full window --- entropy
$\asymp m \log\log m/\log m$ --- admits no bounded budget at all. The impossibility of a short exact
scheme for the primes themselves and the existence of a log-length certified procedure for one sum are
the same pigeonhole read at two entropy scales, which is, we believe, the honest rigorous shape of the
no-closed-form folklore: no more, and no less.

Three questions remain open. (1) \emph{The gap.} Does any scheme, uniform over prime configurations,
decode the window from $(1 + o(1)) \log_2\binom{m}{r}$ bits? We conjecture that smooth
polynomial-weight moments cannot close the $\log m / \log\log m$ factor, and we know no encoding that
does. (2) \emph{Adaptivity.} Extending Theorem~\ref{thm:lb} to schemes that adaptively query real
functionals at bounded precision is routine by the same counting applied to the transcript; a version
with \emph{analog} noise models may be less routine. (3) \emph{Time.} The floor prices
distinguishability, not work: whether the primes of $W$ can be \emph{found} substantially faster than
the elementary square-root barrier is the standing problem of \cite{Polymath}, and nothing here
touches it. The three regimes --- one exact real moment (information complete, extraction hard); many
finite-precision moments (constructive, overpaying); too few bits (impossible) --- delimit what the
moment channels of \cite{CompanionI, CompanionII} can and cannot know, and close the circle this
series of notes set out to draw.

\section*{Disclosure of AI use}

In accordance with publisher policies on the use of artificial intelligence, the author discloses the
following. This paper was developed in an extensive collaboration with Claude (Anthropic; Claude Fable
5, accessed through the claude.ai interface, 2026). The direction of the work --- the founding research
programme of exact moment systems for interval primes, of which this paper is the capstone --- together
with all decisions about scope, claims, and submission, are the author's. Within that direction, the
AI system contributed substantially to the mathematical development and proofs, wrote the verification
code, drafted and revised the manuscript text, and assisted with literature search; the reason for its
use was to enable a single non-specialist author to develop, check, and document the material to a
publishable standard. Every numerical and mathematical claim was subsequently machine-verified by
the recovery experiments reported in Section~6, with the verification code openly available in the cited repository \cite{Code}. The author has
reviewed the entire manuscript and takes full responsibility for its content, including the accuracy
of the mathematics and of all references.

\begin{thebibliography}{99}
\bibitem{CompanionI} C. Gribble, Prime counts and prime sums from multiplication tables: a
self-similar formula system, companion manuscript (2026). % TODO: replace with arXiv ID once posted.
\bibitem{CompanionII} C. Gribble, Certified exact sums of primes over dyadic intervals via a ladder
of generalized factorials, companion manuscript (2026). % TODO: replace with arXiv ID once posted.
\bibitem{Linnik} Yu.~V. Linnik, \emph{The Dispersion Method in Binary Additive Problems}, Amer. Math.
Soc., Providence, 1963.
\bibitem{McKenzie} N. McKenzie, Computing the prime counting function with Linnik's identity,
manuscript (2011), available at \texttt{https://www.primecounting.com}.
\bibitem{Cheb} P.~L. Chebyshev, M\'emoire sur les nombres premiers, \emph{J. Math. Pures Appl.}
\textbf{17} (1852), 366--390.
\bibitem{Macdonald} I.~G. Macdonald, \emph{Symmetric Functions and Hall Polynomials}, 2nd ed., Oxford
Univ. Press, 1995.
\bibitem{Prony} R. de Prony, Essai exp\'erimental et analytique, \emph{J. \'Ecole Polytechnique}
\textbf{1} (1795), 24--76.
\bibitem{Mills} W.~H. Mills, A prime-representing function, \emph{Bull. Amer. Math. Soc.} \textbf{53}
(1947), 604.
\bibitem{Willans} C.~P. Willans, On formulae for the $n$th prime number, \emph{Math. Gaz.} \textbf{48}
(1964), 413--415.
\bibitem{Erdos} P. Erd\H{o}s, Beweis eines Satzes von Tschebyschef, \emph{Acta Litt. Sci. Szeged}
\textbf{5} (1932), 194--198.
\bibitem{Polymath} D.~H.~J. Polymath, Deterministic methods to find primes, \emph{Math. Comp.}
\textbf{81} (2012); arXiv:1009.3956.
\bibitem{Code} C. Gribble, Reference implementations accompanying this paper and its companions,
code repository (2026), \texttt{https://github.com/carlgribble-caa/prime-moments}.
\end{thebibliography}

\end{document}
